
\begin{frame}{Group Fairness Is Not the Only Framing}

    \begin{columns}[T]
        \begin{column}{0.48\textwidth}
            \textbf{Individual fairness}
            \vspace{0.3em}

            ``Similar individuals should receive similar predictions'' (Dwork et al., 2012).

            \vspace{0.5em}
            \begin{itemize}\small
                \setlength\itemsep{0.3em}
                \item Attractive, and it catches cases group parity misses.
                \item Requires a task-specific \textbf{similarity metric} --- which is the whole
                      problem, moved somewhere else.
            \end{itemize}
        \end{column}
        \begin{column}{0.48\textwidth}
            \textbf{Causal / counterfactual fairness}
            \vspace{0.3em}

            ``Would the decision change if this person had belonged to another group, holding
            everything else fixed?'' (Kusner et al., 2017).

            \vspace{0.5em}
            \begin{itemize}\small
                \setlength\itemsep{0.3em}
                \item Directly expresses what most people mean by discrimination.
                \item Needs a causal graph you are willing to defend, and the attribute is often
                      not manipulable in the required sense.
            \end{itemize}
        \end{column}
    \end{columns}

    \vspace{1em}
    \begin{block}{}
        \centering
        \small For an introductory course: use group metrics because you can compute them, and
        remember that they are a screening tool, not a certificate.
    \end{block}
\end{frame}
